
\documentclass{resume} % Use the custom resume.cls style

% Document margins
\usepackage[left=0.75in,top=0.6in,right=0.75in,bottom=0.6in]{geometry}

% Color and hyperlink packages
\usepackage{xcolor}
\usepackage{hyperref}

% Footnote and margin adjustment packages
\usepackage{footnote}
\usepackage{changepage}

% Fontawesome package for icons
\usepackage{fontawesome}

% Tabularx package for custom tables
\usepackage{tabularx}

% Define navyblue color
\definecolor{navyblue}{RGB}{0,54,123}

%----------------------------------------------------------------------------------------
%   Customizations
%----------------------------------------------------------------------------------------

% Define italicitem, bolditem, and plainitem commands
\newcommand{\italicitem}[1]{\item{\textit{#1}}}
\newcommand{\bolditem}[1]{\item{\textbf{#1}}}
\newcommand{\plainitem}[1]{\item{#1}}

% Define user-friendly link command for hyperlinks
\newcommand{\link}[2]{{\href{#1}{#2}}}

\newcommand{\entry}[2]{#1 & #2 \tabularnewline} % Defines an entry with two arguments: #1 for the first column and #2 for the second column

%----------------------------------------------------------------------------------------
%   Define envsection command for defining a new environment section
%----------------------------------------------------------------------------------------

\newcommand{\tableEnv}[2]{%
  \begin{rSection}{#1} % Begin rSection with the given name
    \begin{adjustwidth}{0.0in}{0.1in} % Set the left and right margins
      \begin{tabularx}{\linewidth}{@{} >{\bfseries}l @{\hspace{6ex}} X @{}}
        #2 % Print the content inside the tabularx environment
      \end{tabularx}
    \end{adjustwidth}
  \end{rSection}
}

%----------------------------------------------------------------------------------------
%   Begin document
%----------------------------------------------------------------------------------------

% Set name with navyblue color
\name{\color{navyblue} Luong Tran}

\begin{document}

\printPersonalInfo{
  \personalInfo{
    \tag{Residence}\info{Hanoi, Vietnam}
    \infoSeparator\tag{Telephone number}\info{+84 971-161-526 }
   
    }
  \personalInfo{\tag{E-mail}\info{
  % trankhanhluong432@gmail.com
  trankhanhluong432@gmail.com} \infoSeparator\tag{LinkedIn}\info{\url{www.linkedin.com/in/luongtk43}}}
}

%----------------------------------------------------------------------------------------
%   Objective section
%----------------------------------------------------------------------------------------
% \begin{rSection}{Objective}
%     \begin{itemize}
%         % \item Improve knowledge, work in a professional research environment. 
%         % \item Have several scientific papers published to apply for a scholarship to study abroad.
%         \item Gain valuable knowledge and experience in a professional research environment with the goal of publishing multiple scientific papers to support scholarship applications for studying abroad.
%     \end{itemize}
% \end{rSection}
%----------------------------------------------------------------------------------------

%   Education section
%----------------------------------------------------------------------------------------

\begin{rSection}{Research Interests}

My research interests focus on the effectiveness of Multi-modal Representation Learning and its applications, particularly in Weakly Supervised Learning scenarios. I am also interested in Multi-modal Generative AI and in scaling such systems effectively using Mixture-of-Experts architectures. Recently, I have been exploring vision-based World Models combined with Reinforcement Learning as a core component for AI-driven robotics.
    
\end{rSection}

\begin{rSection}{Education}
    
    % Bachelor's degree entry
    \begin{rSubsectionNoBullet}{\bf Bachelor's Degree in Data Science (Advanced Program)}{Hanoi University of Science and Technology}{Major: Data Science and Artificial Intelligence}{2020 - 2024}
        \italicitem{GPA: 3.75/4.0 (9.38/10)}
        \italicitem{Relevant Coursework: Machine Learning, Data Structures \& Algorithms, Linear Algebra, Probability and Statistics}
    \end{rSubsectionNoBullet}

\end{rSection}

\begin{rSection}{Publication}

% \begin{rSubsectionNoBullet}{\href{https://openaccess.thecvf.com/content/ICCV2025/papers/Tran_More_Reliable_Pseudo-labels_Better_Performance_A_Generalized_Approach_to_Single_ICCV_2025_paper.pdf}{More Reliable Pseudo-labels, Better Performance: A Generalized Approach to Single Positive Multi-label Learning}}{July 2025}{}{} 
% \item
%     \begin{itemize}
%         \item Conference: International Conference on Computer Vision (ICCV 2025)
%         \item Contributions: First author; develop a simple yet effective framework for dynamic pseudo-label generation during training, complemented by a novel loss function designed to robustly learn from diverse pseudo-labels while reducing noise.
%         \item Code: \href{https://github.com/Fsoft-AIC/AEVLP}{github.com/Fsoft-AIC/AEVLP}; \href{https://arxiv.org/abs/2508.20381}{link pre-print}.
%     \end{itemize}
% \end{rSubsectionNoBullet}


% \begin{rSubsectionNoBullet}{\href{https://dl.acm.org/doi/abs/10.1145/3628797.3629002}{Improving Single Positive Multi-label Classification via Knowledge-based Label-weighted Large Loss Rejection}}{December 2023}{}{} 
% \item
%     \begin{itemize}
%         \item Conference: 12th International Symposium on Information and Communication Technology (SOICT 2023)
%         \item Contributions: Co-first author; proposed a knowledge-based label-weighted loss rejection method that significantly improves single-positive multi-label classification, achieving superior performance over state-of-the-art baselines.
%         \item \href{https://dl.acm.org/doi/abs/10.1145/3628797.3629002}{Link paper.}
%     \end{itemize}
% \end{rSubsectionNoBullet}

\begin{rSubsectionPaper}
    {CLIP-FMoE: Scalable CLIP via Fused Mixture-of-Experts with Enforced Specialization, \href{https://openreview.net/forum?id=RErUA45pI9}{\textcolor{cyan}{Link}}}
    {\textbf{Luong Tran}, Lan-Cuong Nguyen, Huynh Dang Nguyen, Dat Nguyen Cong, Dung D. Le, Van Nguyen}
    {The International Conference on Learning Representations, ICLR 2026}
    {}
    
\end{rSubsectionPaper}

\vspace{-0.4cm}

\begin{rSubsectionPaper}
    {More Reliable Pseudo-labels, Better Performance: A Generalized Approach to Single Positive Multi-label Learning, \href{https://openaccess.thecvf.com/content/ICCV2025/papers/Tran_More_Reliable_Pseudo-labels_Better_Performance_A_Generalized_Approach_to_Single_ICCV_2025_paper.pdf}{\textcolor{cyan}{Link}}}
    {\textbf{Luong Tran}, Thieu Vo, Anh Nguyen, Sang Dinh Viet, Van Nguyen}
    {The International Conference on Computer Vision, ICCV 2025}
    {}
    
  \end{rSubsectionPaper}

\vspace{-0.4cm}


\begin{rSubsectionPaper}
    {\bf{Improving Single Positive Multi-label Classification via Knowledge-based Label-weighted Large Loss Rejection}, \href{https://dl.acm.org/doi/abs/10.1145/3628797.3629002}{\textcolor{cyan}{Link}}}
    {Duc Nguyen Quang*, \textbf{Luong Tran}*, Duc Le Hong, Hoan Nguyen Huy, Sang Dinh Viet}
    {The International Symposium on Information and Communication Technology, SoICT 2023}
    {}
    
  \end{rSubsectionPaper}

\vspace{-0.4cm}


% \begin{rSubsectionNoBullet}{\href{https://drive.google.com/file/d/1UiTmF-saMnY7bkJch_Wdd9kJI1MIasSC/view?usp=sharing}{CLIP-FMoE: Scalable CLIP via Fused Mixture-of-Experts
% with Enforced Specialization}}{August 2025}{}{}
% \item
%     \begin{itemize}
%         \item Scale CLIP parameters using a Mixture-of-Experts approach, enabling efficient training with fewer than 3M image–caption pairs (less than 1\% of the original 400M pretraining data).
%         \item Achieve strong improvements in image–text retrieval on both standard and long, complex captions, while maintaining zero-shot classification performance.
%         \item First author; submit to ICLR 2026 (under review); \href{https://drive.google.com/file/d/1UiTmF-saMnY7bkJch_Wdd9kJI1MIasSC/view?usp=sharing}{link pre-print}.
%     \end{itemize}
% \end{rSubsectionNoBullet}



% \begin{rSubsectionNoBullet}{Project in Multimodal Models}{August - October 2024}{LiBMoE}{}
% \item
%     \begin{itemize}
%         \item  Reimplement the HyperRouter algorithm, Mixture of Experts technique, and integrate it into the LLaVA-style models of the LiBMoE framework.
%         \item Conduct training and evaluation experiments, document results, and determine optimal HyperRouter configurations while maintaining fair comparisons with other algorithms.
%         \item Paper: Co-author;  \href{https://arxiv.org/abs/2411.00918}{LIBMoE: A Library for Comprehensive Benchmarking of Mixture of Experts in Large Language Models}
%         \item Code: \href{https://github.com/Fsoft-AIC/LibMoE}{github.com/Fsoft-AIC/LibMoE}.
%     \end{itemize}
    
% \end{rSubsectionNoBullet}

\end{rSection}


\begin{rSection}{Preprint}

\begin{rSubsectionPaper}
    {Precise Video-to-Audio Generation with Cross-Modal Alignment in Latent Space}
    {Thanh V. T. Tran, Ngoc-Son Nguyen, \textbf{Luong Tran}, Long-Khanh Pham, Paarth Neekhara, Shehzeen Samarah Hussain, Van Nguyen}
    {Under review at the IEEE/CVF Conference on Computer Vision and Pattern Recognition, CVPR 2026}
    {}
    
\end{rSubsectionPaper}

\vspace{-0.4cm}


\begin{rSubsectionPaper}
    {LIBMoE: A Library for comprehensive benchmarking Mixture of Experts in Large Language Models, \href{https://arxiv.org/abs/2411.00918}{\textcolor{cyan}{Link}}}
    {Nam V. Nguyen, Thong T. Doan, \textbf{Luong Tran}, Van Nguyen, Quang Pham}
    {Under review at Transactions on Machine Learning Research, TMLR}
    {}
    
\end{rSubsectionPaper}

\vspace{-0.4cm}




\end{rSection}




%   Experience section
%----------------------------------------------------------------------------------------

\begin{rSection}{Research Experience}
    % First work experience entry
    \begin{rSubsectionNoBullet}{AI Center, FPT Software}{July 2024 - present}{AI Resident}{}
    \item Advisor: Dr.\ \textbf{Van Nguyen}, Associate Prof. \textbf{Anh Nguyen}
    \begin{itemize}
        % \item {Design and deploy a multi-agent personalized chatbot aimed at reducing stress and harmful moods among developers, with an MVP released to sales teams for evaluation.}
        % \item {Explore the application of Mixture of Experts to design specialized architectures for Multimodal LLMs, enabling scalable and efficient training through up-cycling without retraining from scratch.}
        % \item {Study Multimodal Alignment Learning for vision–language tasks, with a focus on effective post-pre-training methods for CLIP.}
        \item {Design and deploy a multi-agent personalized chatbot to reduce stress and harmful moods among developers, with an MVP released to sales teams for evaluation.}
        \item {Explore Mixture-of-Experts architectures for Multimodal LLMs to design specialized components that enable scalable and efficient training via up-cycling instead of retraining from scratch.}
        \item {Research about multimodal alignment learning for vision–language tasks, with a focus on effective post-pre-training adaptation methods for CLIP.}
    
    \end{itemize}
    \end{rSubsectionNoBullet}

     % Second work experience entry
    \begin{rSubsectionNoBullet}{Vingroup}{July 2023 - May 2024}{Vingroup AI Engineer Training Program}{}
        % \item Completed a 4-month intensive training program, gaining proficiency in advanced AI concepts, frameworks, and tools, and subsequently 
        \item
            \begin{itemize}
                \item {Complete a 4-month intensive, R\&D-oriented training program, gaining proficiency in advanced Machine Learning, particularly in Computer Vision and Natural Language Processing, and rank in the top 1\% of program participants.}
                % \item {Work with 3 partners to train and fine-tune a Large Language Model for Vietnamese Poem Generation with Pytorch frameworks, resulting model demonstrates the ability to generate poetry across various genres.}
                \item {Design, implement, and deploy lightweight models for vehicle Re-Identification and Face Recognition, empirically reducing error rates to under 2\% on the company’s large-scale urban dataset.}
                \item {Apply and evaluate quantization techniques for YOLOv5 and integrated DeepSort for camera-based vehicle tracking in parking areas, achieving a 25\% reduction in system latency and improving overall deployment efficiency.}
            \end{itemize}
    \end{rSubsectionNoBullet}

    % % First work experience entry
    \begin{rSubsectionNoBullet}{Computer Vision Research Team, BKAI HUST}{October 2020 - July 2024}{Research Assistant}{}
        % \item Review Images Classification on Moods: Employ cutting-edge Vision-Language models to optimize multi-label classification performance using data sourced from the Naver social network..
        \item Advisor: Assistant Prof.\ \textbf{Sang Dinh Viet}
            \begin{itemize}
                \item {Conduct research on projects funded by Naver Corporation.}
                \item {Develop multi-label and semi-supervised learning approaches to build an image tagging system for Naver social network data, addressing challenges of limited and noisy supervision.}
                % \item {Conducting research on real-world problems, particularly focusing on object re-identification. }
            \end{itemize}
    \end{rSubsectionNoBullet}
    
\end{rSection}

%   Project section
%----------------------------------------------------------------------------------------

% \begin{rSection}{Project}

% \begin{rSubsectionNoBullet}{Project in Computer Vision}{September 2024 - March 2025}{Multi-label Learning with Single Positive Labels}{}
% \item
%     \begin{itemize}
%         \item Project on training models for multi-label image classification using only one label per image.
%         \item  Develop a novel framework that leverages vision-language models to augment supervision for multi-label classification with limited annotations.
%         \item Achieve SOTA results by utilizing adaptive pseudo-labeling techniques.
%     \end{itemize}
    
% \end{rSubsectionNoBullet}


% \begin{rSubsectionNoBullet}{Project in Multimodal Models}{August - October 2024}{LiBMoE}{}
% \item
%     \begin{itemize}
%         \item  Reimplement the HyperRouter algorithm, Mixture of Experts technique, and integrate it into the LLaVA-style models of the LiBMoE framework.
%         \item Conduct training and evaluation experiments, document results, and determine optimal HyperRouter configurations while maintaining fair comparisons with other algorithms.
%         \item Paper: Co-author;  \href{https://arxiv.org/abs/2411.00918}{LIBMoE: A Library for Comprehensive Benchmarking of Mixture of Experts in Large Language Models}
%         \item Code: \href{https://github.com/Fsoft-AIC/LibMoE}{github.com/Fsoft-AIC/LibMoE}.
%     \end{itemize}
    
% \end{rSubsectionNoBullet}

% \begin{rSubsectionNoBullet}{Project in LLM Application}{July - October 2024}{Voicebot Receptionist}{}
% \item
%     \begin{itemize}
%         \item Deliver an MVP voicebot receptionist enabling seamless speech-based interaction to handle basic queries about FPT corporate offices.
%         \item Employ prompt engineering techniques to improve response accuracy and user experience, conducted systematic testing, and provided iterative feedback to refine the system.
%         \item Implement retrieval-augmented generation (RAG) techniques to accurately retrieve and present information about FPT.
%     \end{itemize}
% \end{rSubsectionNoBullet}


% \begin{rSubsectionNoBullet}{Project in Natural Language Processing}{January 2024}{Developing a News Classification Model Based on Titles}{}
% \item
%     \begin{itemize}
%         \item Using the LoRA technique (low-rank adaptation) to train a model for categorizing news titles into specific categories, optimizing computational resources for efficiency.
%         \item Implement a FastAPI deployment endpoint to facilitate real-time classification of news titles, ensuring seamless integration and accessibility.
%     \end{itemize}
    
% \end{rSubsectionNoBullet}


% Second project entry
% \begin{rSubsectionNoBullet}{Project in Natural Language Processing}{May - July 2023}{Developing a Healthcare Chatbot specialized in various aspects of common diseases.}{}
% \item
%     \begin{itemize}
%         \item Crawl raw medical information data from Wikipedia and apply the large language model GPT-3.5 to generate a comprehensive question-answer dialogue dataset on common diseases.
%         \item Train and fine-tune the chatbot with the GPT-2 architecture using an artificially self-created dataset on a distributed computing system to establish a coherent and natural question-answering system.
%         \item Improve chatbot reliability and accuracy with the RAG technique, leveraging the Qdrant vector database and a free embedding model from HuggingFace.
%         % \item{Implement a user-friendly interface for the chatbot to interact with users via text commands.}
%     \end{itemize}
    
% \end{rSubsectionNoBullet}

% % First project entry
% \begin{rSubsectionNoBullet}{Project in Data Science}{November - December 2022}{Lyric-based music genre classification}{}  
% \item
%     \begin{itemize}
%         \item {Perform ETL processes to analyze the data crawled from Spotify, which is song lyrics, for the classification task.}
%         \item {Implement suitable Transformer models to solve the problem based on the characteristics of the data and fine-tune the models to obtain the best result.}
%     \end{itemize}
    
% \end{rSubsectionNoBullet}



% \end{rSection}





%----------------------------------------------------------------------------------------

%----------------------------------------------------------------------------------------



%----------------------------------------------------------------------------------------
%   Activities section
%----------------------------------------------------------------------------------------

\begin{rSection}{Awards \& Honors}

\begin{rSubsectionNoBullet}{Best Presentation Award}{July 2024}{Hanoi University of Science and Technology (HUST)}{}
    \plainitem{Awarded for outstanding thesis presentation in the 2024.2 Computer Vision Council graduation.}
\end{rSubsectionNoBullet}
% First activity entry
\begin{rSubsectionNoBullet}{JUNCTIONX HANOI 2023}{April 2023}{Co-organized by Junction Vietnam and Viettel Group}{}
\item{
    \begin{itemize}
        \item Track Winner of Cloud Track: Cloud Application for Community
         \item{Develop the EcoFrenzy app, a social media platform that challenges its users to take green and fun actions in their everyday lives, with the goal of reducing air pollution caused by transportation.}
    \end{itemize}
}
   
\end{rSubsectionNoBullet}

\end{rSection}


%----------------------------------------------------------------------------------------
% Technical skills section
%----------------------------------------------------------------------------------------


\tableEnv{Skills}{
    \entry{Programming Languages/Frameworks}{Python, PyTorch, C/C++, Java, LangChain, OpenCV}
    \entry{Languages}{English (7.0 IELTS)}
}

\begin{rSection}{References}
    \begin{itemize}{}{\itemsep=0.4em}  % no bullets, vertical spacing
    
        \item \begin{rReference}
            {\href{https://cgi.csc.liv.ac.uk/~anguyen/}{\textcolor{cyan}{Associate Prof.\ Anh Nguyen}}}
            {Department of Computer Science, University of Liverpool}
            {Email: anh.nguyen@liverpool.ac.uk}
            {}
        \end{rReference}

        \item \begin{rReference}
            {\href{https://www.linkedin.com/in/vannguyen8/}{\textcolor{cyan}{Dr.\ Van Nguyen}}}
            {AI Transformation Lab Manager, Head of Research, FPT Software}
            {Email: vannth19@fpt.com}
            {}
        \end{rReference}

        \item \begin{rReference}
            {\href{https://users.soict.hust.edu.vn/sangdv/}{\textcolor{cyan}{Assistant Prof.\ Sang Dinh Viet}}}
            {Vice Dean, School of Information and Communication Technology, HUST}
            {Email: sangdv@soict.hust.edu.vn}
            {}
        \end{rReference}

    \end{itemize}
\end{rSection}



\end{document}